\section{Descrizione del Sistema}

L'applicazione presa in esame in questo progetto è una \textbf{web application bancaria} che fornisce funzionalità di home banking e gestione dei servizi bancari di base. L'applicazione si propone di coprire tre principali aree di utilizzo, ciascuna destinata a una specifica categoria di utenti:

\subsection{Categorie di Utenti}

\subsubsection{Clienti}
I clienti della banca possono accedere a un'interfaccia intuitiva per gestire le principali operazioni bancarie online, come:
\begin{itemize}
	\item Visualizzazione del saldo e della lista movimenti;
	\item Bonifici e trasferimenti di denaro;
	\item Gestione dei propri dati anagrafici.
\end{itemize}

\subsubsection{Employee}
I dipendenti della banca possono accedere a un pannello dedicato che consente di:
\begin{itemize}
	\item Visualizzare e modificare i dati dei clienti;
	\item Gestire operazioni richieste dai clienti in filiale;
	\item Gestire le credenziali per l'accesso dei clienti.
\end{itemize}

\subsubsection{Admin}
Gli admin hanno accesso completo alla gestione dell'utenza e delle funzionalità del sistema, con la possibilità di:
\begin{itemize}
	\item Creare, modificare e cancellare account (clienti e employee);
	\item Gestire le assegnazioni fra clienti e employee;
	\item Monitorare le attività del sistema.
\end{itemize}

\newpage 

\subsection{Principali Schermate dell'Applicazione}
\subsubsection{Schermate utente ancora non autenticato}

La Homepage presenta un bottone in alto a sinistra che permette all'utente di accedere. 

\begin{figure}[htbp]
	\centering
	\includegraphics[width=0.8\textwidth]{schermata_principale}
	\caption{Homepage}
	\label{fig:schermata_principale}
\end{figure}

Una volta cliccato il bottone di login, l'utente vedrà due opzioni:
\begin{itemize}
	\item Accesso come cliente;
	\item Accesso come lavoratore (ovvero come employee oppure admin);
\end{itemize} 

\begin{figure}[htbp]
	\centering
	\includegraphics[width=0.8\textwidth]{selezione_tipologia_accesso}
	\caption{Schermata per selezionare la tipologia di accesso}
	\label{fig:selezione_tipologia_accesso}
\end{figure}

\clearpage

Nella pagina di login l'utente potrà inserire username e password. 

\begin{figure}[htbp]
	\centering
	\includegraphics[width=0.8\textwidth]{login_cliente}
	\caption{Schermata login}
	\label{fig:login_cliente}
\end{figure}

\begin{figure}[htbp]
	\centering
	\includegraphics[width=0.5\textwidth]{login_compilato_cliente}
	\caption{Schermata login compilata}
	\label{fig:login_compilato_cliente}
\end{figure}